\documentclass[11pt]{article}

\usepackage[margin=1.02in]{geometry}
\usepackage{amsmath,amssymb,amsthm,mathtools}
\usepackage{array}
\usepackage{booktabs}
\usepackage{enumitem}
\usepackage{microtype}
\usepackage{xcolor}
\usepackage[hidelinks]{hyperref}
\usepackage{fancyhdr}

\definecolor{deepblue}{HTML}{173B57}
\hypersetup{
  colorlinks=true,
  linkcolor=deepblue,
  citecolor=deepblue,
  urlcolor=deepblue,
  pdftitle={Tensor-local constraints in the exceptional unitary Hecke Yang--Baxter class},
  pdfauthor={Alec Kriebel}
}

\pagestyle{fancy}
\fancyhf{}
\fancyhead[L]{\small\textsc{Exceptional Hecke constraints}}
\fancyhead[R]{\small\textsc{A. Kriebel}}
\fancyfoot[C]{\thepage}
\setlength{\headheight}{14pt}

\newtheorem{theorem}{Theorem}[section]
\newtheorem{lemma}[theorem]{Lemma}
\newtheorem{proposition}[theorem]{Proposition}
\newtheorem{corollary}[theorem]{Corollary}
\theoremstyle{definition}
\newtheorem{definition}[theorem]{Definition}
\theoremstyle{remark}
\newtheorem{remark}[theorem]{Remark}
\newtheorem{problem}[theorem]{Open problem}

\newcommand{\Id}{\mathbf 1}
\newcommand{\Tr}{\operatorname{Tr}}
\newcommand{\End}{\operatorname{End}}
\newcommand{\rank}{\operatorname{rank}}
\newcommand{\OSR}{\operatorname{OSR}}
\newcommand{\Ann}{\operatorname{Ann}}
\newcommand{\bbC}{\mathbb C}
\newcommand{\cA}{\mathcal A}
\newcommand{\cB}{\mathcal B}
\newcommand{\cC}{\mathcal C}
\newcommand{\cX}{\mathcal X}
\newcommand{\cY}{\mathcal Y}
\newcommand{\cO}{\mathcal O}
\newcommand{\cE}{\mathscr E}
\newcommand{\cS}{\mathfrak S}
\newcommand{\HS}{\mathrm{HS}}

\title{\vspace{-1.4cm}\textbf{Tensor-local constraints in the exceptional}\\[0.2em]
\textbf{unitary Hecke Yang--Baxter class}}
\author{Alec Kriebel\thanks{This work was developed with substantial
assistance from OpenAI ChatGPT. The exact scope of that assistance is
disclosed at the end of the paper.}}
\date{Version 0.1.0 -- 29 July 2026}

\begin{document}
\maketitle

\begin{abstract}
Let \(q=e^{i\pi/3}\). The only existence family left unresolved by the
classification of unitary two-eigenvalue Yang--Baxter matrices is
\[
  \left[q,\frac12,d\right],\qquad d>2\ \text{even}.
\]
An explicit solution is known in local dimension four and hence in every
dimension divisible by four; the dimensions \(d\equiv2\pmod4\) remain open.
We determine the exact structural frontier of this problem.

First, every matrix in the exceptional class is automatically standard:
both partial traces of its negative spectral projection equal
\((d/2)\Id_d\). Its tensor-power representation therefore has the
\(\eta=1/2\) Markov trace and is faithful after passage to the
Jones--Wenzl trace quotient \(H_n(3,6)\). We compute every simple
multiplicity in this representation:
\[
  m_{\lambda,n}=D_\lambda\left(\frac d2\right)^n,
  \qquad D_\lambda\in\{1,2,3\}.
\]
Consequently the complete abstract Hecke-tower arithmetic permits every
even \(d\); a four-divisibility obstruction must use strict tensor locality.

We then prove that an exceptional reflection of operator-Schmidt rank at
most three exists exactly when \(4\mid d\). At rank four, projecting the
full cubic relation modulo the intrinsic Schmidt supports yields two
joint-sandwich identities. If either associated sandwich map is injective,
then \(4\mid d\); hence every rank-four candidate in
\(d\equiv2\pmod4\) has nonzero Hermitian traceless common annihilators on
both tensor legs.

Finally, any hypothetical dimension-six solution is nonrestrictable,
satisfies
\[
  \cC_L(P)\cap\cC_R(P)=\bbC\Id_6,
\]
and lies outside the primitive-Weyl Bell-diagonal and four-product
Clifford-frame mechanisms. A four-dimensional one-sided square restriction
is not excluded. Thus the complete spectrum remains open, but any
dimension-six witness must be a high-Schmidt-rank, genuinely tensor-local
object outside the standard construction mechanisms isolated here.
\end{abstract}

\medskip
\noindent\textbf{Status.}
This is an unrefereed, AI-assisted research note. The proofs, exact
verifiers, adversarial audits, and computational provenance are public.
Priority claims are provisional and no specialist peer review is claimed.

\section{Introduction and main results}\label{sec:intro}

Let \(V\cong\bbC^d\). A unitary
\(R\in\End(V\otimes V)\) is a Yang--Baxter operator when
\begin{equation}
  R_{12}R_{23}R_{12}=R_{23}R_{12}R_{23}
  \quad\text{on }V^{\otimes3}.
  \label{eq:YBE}
\end{equation}
Lechner~\cite{Lechner2026} classified the possible characters of unitary
Yang--Baxter matrices with two eigenvalues, leaving one potentially empty
family:
\[
  \left[e^{i\pi/3},\frac12,d\right],
  \qquad d>2\text{ even}.
\]
The notation records the nonnegative Hecke phase, the normalized rank of
the \((-1)\)-spectral projection, and the local dimension.

The dimension-two class is empty. In~\cite{Kriebel2026d4} we constructed an
exact \(16\times16\) unitary representative in local dimension four.
Identity stabilization then gives representatives in every dimension
divisible by four. The unresolved part of the spectrum is therefore
\[
  d=6,10,14,\ldots.
  \label{eq:open-dims}
\]
This question is also the strict ordinary-localization problem associated
with the \(\cC(\mathfrak{sl}_3,6)\) Jones--Wenzl braid representations
studied by Galindo, Hong, and Rowell~\cite{GHR2013}.

It is tempting to seek the missing factor of two in the representation
theory of \(H_n(3,6)\). One of the main conclusions of this paper is that
this cannot work at the level of abstract simple multiplicities, central
ranks, or their branching recurrences: all of that arithmetic is compatible
with every even local dimension. The extra restriction, if it exists, must
come from realizing every adjacent Hecke generator by the two shifts of one
common two-site matrix.

We summarize the results in two theorems. Precise definitions appear in
Section~\ref{sec:setup}.

\begin{theorem}[Automatic tower structure and its limitation]
\label{thm:tower-summary}
Let \(P\in\End(V\otimes V)\) be an exceptional projection in local
dimension \(d\). Then:
\begin{enumerate}[label=\textup{(\roman*)}]
\item
\[
  \Tr_1P=\Tr_2P=\frac d2\Id_d.
\]
No finite-dimensional irreducibility assumption is needed.
\item The normalized tensor trace is the \(\eta=1/2\) Markov trace at
every level, and the local Hecke representation induces an injection
\[
  H_n(3,6)\longrightarrow\End(V^{\otimes n}).
\]
Equivalently, the kernel of the raw Hecke representation is exactly the
Markov-trace annihilator.
\item If \(S_{\lambda,n}\) is a simple \(H_n(3,6)\)-module and
\(D_\lambda\in\{1,2,3\}\) is the associated quantum dimension, then its
tensor-space multiplicity is
\[
  m_{\lambda,n}=D_\lambda\left(\frac d2\right)^n.
\]
These multiplicities and all their restriction recurrences are integral
exactly when \(d\) is even.
\end{enumerate}
Thus ordinary Hecke simple-block, central-rank, Markov-trace, and branching
arithmetic cannot force \(4\mid d\).
\end{theorem}

\begin{theorem}[Tensor-local frontier]\label{thm:frontier-summary}
Let \(H=\Id-2P\) be the reflection associated with an exceptional
projection.
\begin{enumerate}[label=\textup{(\roman*)}]
\item If a rank-\(r\) projection belongs to either one-leg commutant of
\(P\), then
\[
  8\mid rd^2.
\]
\item An exceptional reflection with \(\OSR(H)\le3\) exists in dimension
\(d\) if and only if \(4\mid d\). Hence every hypothetical solution in
\(d\equiv2\pmod4\) has \(\OSR(H)\ge4\).
\item Suppose \(\OSR(H)=4\), with intrinsic Schmidt supports
\(\cA,\cB\). If either joint-sandwich map
\[
  \cS_{\cB\mid\bbC\Id+\cA}
  \quad\text{or}\quad
  \cS_{\cA\mid\bbC\Id+\cB}
\]
is injective, then \(4\mid d\). Consequently, in
\(d\equiv2\pmod4\), both maps have nonzero Hermitian traceless
coefficient-matrix annihilators.
\item Every square-invariant local subspace inherits a balanced
exceptional solution and has even dimension. A least-dimensional solution
with \(d\equiv2\pmod4\) is nonrestrictable. In particular, no
dimension-six solution is restrictable.
\item In dimension six,
\[
  \cC_L(P)\cap\cC_R(P)=\bbC\Id_6.
\]
This does not assert that either individual leg commutant is scalar.
\item A primitive-Weyl Bell-diagonal exceptional reflection, and every
reflection in the four-product Clifford-frame class defined in
Section~\ref{sec:model-classes}, has \(4\mid d\).
\end{enumerate}
\end{theorem}

Theorem~\ref{thm:frontier-summary}(iii) is the outcome of a final bounded
attack on unrestricted operator-Schmidt rank four. It uses the full cubic,
not a Clifford or symmetry ansatz. It does not prove that either sandwich
map is always injective.

The exact dimension-six conclusion is therefore narrower than
``tensor-locally irreducible.'' A rank-two square-invariant subspace is
impossible, while a rank-four one-sided square-invariant subspace remains
possible in principle; if it exists, its two-dimensional orthogonal
complement cannot also have invariant tensor square.

\section{Exceptional projections and conventions}\label{sec:setup}

Fix throughout
\[
  q=e^{i\pi/3}.
\]
An \emph{exceptional projection} is an orthogonal projection
\[
  P=P^*=P^2\in\End(V\otimes V),\qquad
  \rank P=\frac{d^2}{2},
  \label{eq:projection}
\]
such that, with \(P_1=P\otimes\Id\) and \(P_2=\Id\otimes P\),
\begin{equation}
  P_1P_2P_1-P_2P_1P_2=\frac13(P_1-P_2).
  \label{eq:cubic-P}
\end{equation}
The associated reflection and unitary Hecke operator are
\begin{equation}
  H=\Id-2P,\qquad
  R=q\Id-(1+q)P
   =\frac{q-1}{2}\Id+\frac{q+1}{2}H.
  \label{eq:H-R}
\end{equation}
A direct expansion shows that~\eqref{eq:cubic-P} is equivalent to
\begin{equation}
  H_1H_2H_1-H_2H_1H_2=\frac13(H_1-H_2),
  \label{eq:cubic-H}
\end{equation}
and that~\eqref{eq:cubic-H} is equivalent to the Yang--Baxter equation
for \(R\). Moreover,
\[
  H=H^*,\quad H^2=\Id,\quad\Tr H=0,
\]
and \(R\) is unitary with eigenvalues \(-1,q\), each of multiplicity
\(d^2/2\).

The operator-Schmidt rank of \(X\in M_d\otimes M_d\), denoted
\(\OSR(X)\), is the least \(r\) for which
\[
  X=\sum_{j=1}^r A_j\otimes B_j.
\]
Our rank statements concern the reflection \(H\), not \(R\). Once
automatic standardness is proved, the Schmidt supports of \(H\) consist of
traceless matrices. Hence the identity term in~\eqref{eq:H-R} adds one
independent Schmidt direction:
\begin{equation}
  \OSR(R)=\OSR(H)+1.
  \label{eq:osr-shift}
\end{equation}

Define the two one-leg commutants
\begin{equation}
\begin{split}
  \cC_L(P)&=\{x\in M_d:[x\otimes\Id,P]=0\},\\
  \cC_R(P)&=\{x\in M_d:[\Id\otimes x,P]=0\}.
\end{split}
\label{eq:leg-commutants}
\end{equation}
These are two subalgebras of the same \(M_d\), after identifying the two
copies of the local space. Their intersection is meaningful, but it must
not be confused with either algebra individually.

\begin{definition}\label{def:restrictable}
A nonzero subspace \(W\subseteq V\) is \emph{square-invariant} when
\[
  R(W\otimes W)\subseteq W\otimes W.
\]
The solution is \emph{restrictable} when there is a proper nonzero
\(W\subset V\) for which both \(W^{\otimes2}\) and
\((W^\perp)^{\otimes2}\) are invariant.
\end{definition}

Because \(R\) is unitary, every invariant tensor square is automatically
reducing. Definition~\ref{def:restrictable} agrees with the use in
Conti--Lechner~\cite[\S6.1]{ContiLechner}.

\section{Automatic standardness and the trace quotient}
\label{sec:standardness}

The first theorem removes a potentially serious hypothesis from the
dimension problem.

\begin{theorem}[Automatic standardness]\label{thm:standardness}
Every exceptional projection satisfies
\[
  \Tr_1P=\Tr_2P=\frac d2\Id_d.
  \label{eq:standardness}
\]
\end{theorem}

\begin{proof}
The Hecke operator \(R\) in~\eqref{eq:H-R} has spectrum
\(\{-1,q\}\). It has no pair of opposite eigenvalues, since that would
require \(q=1\). Lechner's no-opposite-spectrum theorem
\cite[Proposition 2.4 and Lemma 3.1]{Lechner2026} therefore makes the
associated braid character a Markov trace. This derives the needed
subfactor irreducibility; it is not an assumption on the finite matrix
representation.

Lechner's partial-trace criterion
\cite[Proposition 2.3]{Lechner2026} says that the Markov property is
equivalent to
\[
  \varphi(R)=\tau(R)\Id_d,\qquad
  \varphi=\frac1d(\Tr\otimes\operatorname{id}),
\]
where \(\tau\) is normalized matrix trace. Put
\[
  \eta=\tau(P)=\frac{\rank P}{d^2}=\frac12.
\]
Using \(R=q\Id-(1+q)P\) and \(1+q\ne0\), comparison gives
\[
  \varphi(P)=\eta\Id_d=\frac12\Id_d.
\]
This is the first identity in~\eqref{eq:standardness}.

For completeness, apply the same argument to the tensor-flipped operator
\(FRF\). Reversal of the three tensor positions preserves the Yang--Baxter
equation, and \(FRF\) has the same spectrum. Its first partial trace is the
second partial trace of \(R\), proving the second identity. One may instead
use the equality of the two partial traces proved in
\cite[Theorem 5.10]{ContiLechner}.
\end{proof}

\begin{corollary}\label{cor:zero-marginals}
Every exceptional reflection satisfies
\[
  \Tr_1H=\Tr_2H=0.
\]
In every minimal Schmidt decomposition of \(H\), both intrinsic Schmidt
supports consist of traceless matrices.
\end{corollary}

\begin{proof}
The first claim follows from \(H=\Id-2P\). For the second, write
\(H=\sum_jA_j\otimes B_j\) with each family linearly independent.
Then
\(\Tr_2H=\sum_j(\Tr B_j)A_j=0\), so every \(\Tr B_j=0\); the other
leg is identical.
\end{proof}

Let \(H_n(q)\) be the specialized Hecke algebra generated by projections
\(e_1,\ldots,e_{n-1}\), with braid generators
\[
  g_i=q\Id-(1+q)e_i.
\]
The local representation is
\[
  \rho_n:H_n(q)\longrightarrow\End(V^{\otimes n}),
  \qquad e_i\longmapsto P_i.
\]
Write \(\tau_n=d^{-n}\Tr\) for normalized matrix trace.

\begin{proposition}[Faithful trace quotient]\label{prop:trace-quotient}
For every \(n\), \(\tau_n\circ\rho_n\) is the Markov trace
\(\mu_{1/2}\). If
\[
  \Ann_n
  =\{x\in H_n(q):\mu_{1/2}(yx)=0
      \text{ for every }y\in H_n(q)\},
\]
then
\[
  \ker\rho_n=\Ann_n.
\]
Consequently \(\rho_n\) induces an injection
\[
  H_n(3,6)=H_n(q)/\Ann_n
  \hookrightarrow\End(V^{\otimes n}).
\]
\end{proposition}

\begin{proof}
For \(x\in H_n(q)\), Theorem~\ref{thm:standardness} gives
\[
\begin{aligned}
 \tau_{n+1}\!\left((\rho_n(x)\otimes\Id)P_n\right)
 &=d^{-(n+1)}
   \Tr\!\left(\rho_n(x)\Tr_{n+1}(P_n)\right)\\
 &=\frac12\,\tau_n(\rho_n(x)).
\end{aligned}
\]
Thus the Markov property holds at every level; it is not inferred merely
from a two-strand trace coincidence. Uniqueness identifies this trace with
\(\mu_{1/2}\).

Wenzl's positive-trace parameter in the present convention is
\[
  \eta_{\ell,k}
  =
  \frac{\sin(\pi(k-1)/\ell)}
       {2\cos(\pi/\ell)\sin(\pi k/\ell)}.
\]
At \((\ell,k)=(6,3)\), it is \(1/2\), identifying the quotient with
\(H_n(3,6)\); see~\cite[Theorem 3.3]{Lechner2026} and
\cite{Wenzl1988}.

If \(x\in\Ann_n\), choose \(y=x^*\). Then
\[
  0=\mu_{1/2}(x^*x)
  =\tau_n(\rho_n(x)^*\rho_n(x)).
\]
Faithfulness of ordinary matrix trace gives \(\rho_n(x)=0\). The reverse
inclusion is immediate. This proves equality of the kernel and
annihilator. Notice that the raw specialized Hecke representation need
not be faithful; faithfulness is asserted only after passage to the trace
quotient.
\end{proof}

\section{Complete tower multiplicities and their limitation}
\label{sec:multiplicities}

The simple \(H_n(3,6)\)-modules are indexed by admissible Young diagrams.
Equivalently, after subtracting the third row, their labels lie in the
ten-point alcove
\[
  (a,b),\qquad a,b\ge0,\quad a+b\le3.
\]
The associated quantum dimensions are
\begin{equation}
 D_{a,b}=
 \frac{
  \sin((a+1)\pi/6)\sin((b+1)\pi/6)
  \sin((a+b+2)\pi/6)}
 {\sin(\pi/6)^2\sin(2\pi/6)}.
 \label{eq:quantum-dim}
\end{equation}
Their values are
\[
\begin{array}{c|c}
D_{a,b}&(a,b)\\ \hline
1&(0,0),(3,0),(0,3)\\
3&(1,1)\\
2&(1,0),(0,1),(2,0),(0,2),(2,1),(1,2).
\end{array}
\]
Adding a box gives the usual alcove branching, and direct substitution
in~\eqref{eq:quantum-dim} yields
\begin{equation}
  \sum_{\nu:\lambda\nearrow\nu}D_\nu=2D_\lambda.
  \label{eq:PF}
\end{equation}

\begin{theorem}[All simple multiplicities]\label{thm:multiplicities}
Decompose the tensor representation as
\[
  V^{\otimes n}
  \cong
  \bigoplus_\lambda
  S_{\lambda,n}\otimes\bbC^{m_{\lambda,n}}.
\]
Then
\[
  \boxed{
  m_{\lambda,n}
  =D_\lambda\left(\frac d2\right)^n.}
  \label{eq:multiplicity}
\]
\end{theorem}

\begin{proof}
Let \(f_{\lambda,n}=\dim S_{\lambda,n}\). The Wenzl trace of a
minimal projection in the \(\lambda\)-block is
\[
  t_{\lambda,n}=\frac{D_\lambda}{2^n};
\]
equivalently, this is the categorical trace under the identification
\(H_n(3,6)\cong\End_{\cC}(X^{\otimes n})\) with
\(\operatorname{FPdim}(X)=2\). The recursion for these weights is exactly
\eqref{eq:PF} divided by \(2^{n+1}\).

Under \(\rho_n\), a minimal projection in the \(\lambda\)-block has
ordinary matrix rank \(m_{\lambda,n}\), hence normalized trace
\(m_{\lambda,n}/d^n\). Proposition~\ref{prop:trace-quotient} equates this
with \(D_\lambda/2^n\), proving~\eqref{eq:multiplicity}.
\end{proof}

\begin{corollary}[Exact arithmetic limitation]\label{cor:tower-limitation}
All simple multiplicities, central-idempotent ranks, and represented
restriction recurrences are integral if and only if \(d\) is even. In
particular, this complete abstract tower arithmetic cannot distinguish
\(d=6\) from \(d=4\).
\end{corollary}

\begin{proof}
If \(d=2s\), formula~\eqref{eq:multiplicity} is
\(m_{\lambda,n}=D_\lambda s^n\), an integer at every level. Moreover,
\eqref{eq:PF} gives
\[
  \sum_{\nu:\lambda\nearrow\nu}m_{\nu,n+1}
  =d\,m_{\lambda,n},
\]
which is exactly the tensor-space restriction recurrence. If \(d\) is
odd, already the two-strand negative spectral rank \(d^2/2\) is not an
integer.
\end{proof}

\begin{remark}[What the limitation proves]\label{rem:abstract-only}
For \(d=2s\), the abstract modules
\[
  \bigoplus_\lambda
  S_{\lambda,n}\otimes\bbC^{D_\lambda s^n}
\]
have the required dimensions, traces, and restriction multiplicities.
At \(d=6\), the three-strand common-one, common-zero, and generic
two-projection multiplicities are \(27,27,81\). Exact four-strand models
also realize the expected simple corners and pairings with the odd
multiplicity \(81\).

These are abstract Hecke modules, not a Yang--Baxter matrix on
\(\bbC^6\otimes\bbC^6\). They do not realize
\(P_{12}=P\otimes\Id_6\) and \(P_{23}=\Id_6\otimes P\) for one common
two-site projection. The insufficiency conclusion is therefore precisely
that ordinary Hecke block data and their branching cannot prove
four-divisibility; strict same-\(P\) tensor coherence remains available as
an obstruction.
\end{remark}

\section{Two broad construction mechanisms}
\label{sec:model-classes}

We record two invariant model classes which are broad enough to include
many natural roots-of-unity constructions, but which still force
four-divisibility.

\subsection{Primitive-Weyl Bell-diagonal reflections}

Let
\[
 X|j\rangle=|j+1\rangle,\qquad
 Z|j\rangle=\zeta^j|j\rangle,\qquad
 \zeta=e^{2\pi i/d},
\]
and define the generalized Bell basis
\[
 |\Phi_{a,b}\rangle
 =
 \frac1{\sqrt d}\sum_{j\in\mathbb Z_d}
 \zeta^{bj}|j,j+a\rangle.
\]
The pair \(X,Z\) is primitive: \(ZX=\zeta XZ\), and the displayed Bell
lines are the simple joint eigenspaces of \(X\otimes X\) and
\(Z^{-1}\otimes Z\).

\begin{theorem}[Bell-diagonal divisibility]\label{thm:bell}
Let
\[
 H=\sum_{a,b\in\mathbb Z_d}
 h_{a,b}|\Phi_{a,b}\rangle\langle\Phi_{a,b}|,
 \qquad h_{a,b}\in\{\pm1\},\quad
 \sum_{a,b}h_{a,b}=0.
\]
If \(H\) satisfies the exceptional cubic, then \(4\mid d\).
\end{theorem}

\begin{proof}
Each Bell line is maximally entangled, so the balance condition gives
\(\Tr_1H=\Tr_2H=0\).
Put
\[
 S=H_{12},\qquad T=H_{23},\qquad U=ST.
\]
The operator \(U\) commutes with
\[
 \mathsf X=X_1X_2X_3,\qquad
 \mathsf Z=Z_1^{-1}Z_2Z_3^{-1}.
\]
These obey
\[
  \mathsf X\mathsf Z=\zeta\,\mathsf Z\mathsf X,
  \qquad \mathsf X^d=\mathsf Z^d=\Id.
\]
Primitivity implies that their generated algebra is isomorphic to
\(M_d(\bbC)\). One direct check is that the \(d^2\) monomials
\(\mathsf X^r\mathsf Z^s\) are trace orthogonal. Its representation on
the \(d^3\)-dimensional three-site space is therefore a direct sum of
\(d^2\) copies of the unique \(d\)-dimensional simple module. Since \(U\)
commutes with this algebra, every eigenvalue multiplicity of \(U\) is
divisible by \(d\).

Using \(S^2=T^2=\Id\), the exceptional cubic gives
\[
 (U-\Id)\left(U^2+\frac23U+\Id\right)=0.
\]
The roots are
\[
 1,\qquad
 \lambda_\pm=-\frac13\pm\frac{2\sqrt2}{3}i.
\]
Conjugation by \(S\) sends \(U\) to \(U^{-1}\), so the
\(\lambda_+\)- and \(\lambda_-\)-multiplicities agree. Call the three
multiplicities \(m_1,m,m\). The zero marginals give
\[
 \Tr(U)
 =\Tr(H_{12}H_{23})=0.
\]
Since \(\lambda_++\lambda_-=-2/3\),
\[
  m_1-\frac23m=0,\qquad m_1+2m=d^3.
\]
Hence
\[
  m_1=\frac{d^3}{4},\qquad m=\frac{3d^3}{8}.
\]
Divisibility of \(m\) by \(d\) gives \(8\mid3d^2\), and therefore
\(4\mid d\).
\end{proof}

\begin{corollary}
The complete primitive-Weyl Bell-diagonal branch is empty in dimension
six.
\end{corollary}

\begin{remark}
Primitivity is essential: it is what makes the common three-site Weyl
algebra a full matrix algebra. The theorem does not assert that an
arbitrary exceptional reflection admits Bell stabilizers. A separate exact
enumeration shows that all \(12{,}870\) balanced sign tables also fail in
the fixed \(d=4\) Bell basis; that finite calibration is recorded in the
supplement rather than used here.
\end{remark}

\subsection{Four-product Clifford frames}

\begin{theorem}[Clifford-frame divisibility]\label{thm:clifford}
Let \(d\) be even and suppose
\[
 H=\sum_{j=1}^4c_jA_j\otimes B_j,
 \qquad c_j\in\mathbb R\setminus\{0\},
\]
where every \(A_j,B_j\) is a traceless Hermitian involution and the four
\emph{product} involutions
\[
 A_j\otimes B_j
\]
anticommute pairwise. Then \(4\mid d\).
\end{theorem}

\begin{proof}
For \(i\ne j\), product anticommutation says
\[
 (A_iA_j)\otimes(B_iB_j)
 =-(A_jA_i)\otimes(B_jB_i).
\]
Equality of nonzero simple tensors gives a scalar \(\epsilon_{ij}\).
Involutivity forces \(\epsilon_{ij}^2=1\). Thus each local pair commutes
or anticommutes, and the choices on the two legs are complementary.

Let \(G_A,G_B\) be the alternating \(4\times4\) binary matrices which
record local anticommutation. If \(K_4\) denotes the adjacency matrix of
the complete four-vertex graph, then
\[
  G_A+G_B=K_4
  \quad\text{over }\mathbb F_2.
\]
If a family of Hermitian involutions has binary commutation matrix of
rank \(2r\), symplectic Gram--Schmidt produces \(r\) mutually commuting
Weyl pairs. They generate \(M_{2^r}(\bbC)\), so \(2^r\mid d\).

Suppose \(d=2s\) with \(s\) odd. Neither \(G_A\) nor \(G_B\) can have
rank four. Since \(K_4\) has binary rank four and alternating ranks are
even, the rank inequality forces
\[
 \rank_{\mathbb F_2}G_A
 =\rank_{\mathbb F_2}G_B=2.
\]
We use the following elementary four-vertex fact: if an alternating
binary matrix \(G\) and its complement \(K_4+G\) both have rank two,
then one of the two graphs has an isolated vertex. To see this, write the
six edge variables as \(a,b,c,d,e,f\). Rank two gives
\[
 af+be+cd=0,
\]
and the complement gives
\[
 1+(a+b+c+d+e+f)+(af+be+cd)=0.
\]
Thus \(G\) has an odd number of edges. A three-edge graph without an
isolated vertex is a path or a star; the path has nonzero Pfaffian and the
star has an isolated vertex in the complement. A five-edge graph has
isolated vertices in its complement.

Let \(U_0\) be an isolated generator on the resulting leg. The same
rank-two graph contains an anticommuting pair \(C,D\). On
\(\bbC^{2s}\), such a pair is unitarily equivalent to
\[
 C=Z\otimes\Id_s,\qquad D=X\otimes\Id_s.
\]
Since \(U_0\) commutes with both, \(U_0=\Id_2\otimes L\) for a Hermitian
involution \(L\in M_s\). But \(s\) is odd, so
\(\Tr L\ne0\), contradicting the assumed tracelessness of \(U_0\).
\end{proof}

\begin{remark}
The load-bearing hypothesis is pairwise anticommutation of the four
product terms. The theorem does not say that either local family is
pairwise anticommuting, nor does it classify arbitrary
operator-Schmidt-rank-four involutions. Their Schmidt factors need not be
involutions and involutivity may result from cancellations among nonzero
cross anticommutators.
\end{remark}

\section{Square restrictions and the dimension-six frontier}
\label{sec:d6-frontier}

The following four-strand calculation makes balance hereditary under a
square restriction.

\begin{lemma}[The two one-dimensional four-strand traces]
\label{lem:four-strand-traces}
Let \(\mu_\eta\) be the Hecke Markov trace with
\(\mu_\eta(e_i)=\eta\), at \(q=e^{i\pi/3}\). If \(e_+\) and \(e_-\) are
the four-strand \(q\)-symmetrizer and \(q\)-antisymmetrizer, then
\begin{equation}
\begin{split}
 \mu_\eta(e_+)
 &=\frac{(1-\eta)(2-3\eta)(1-2\eta)}2,\\
 \mu_\eta(e_-)
 &=\frac{\eta(3\eta-1)(2\eta-1)}2.
\end{split}
\label{eq:four-strand-traces}
\end{equation}
Their unique common zero is \(\eta=1/2\).
\end{lemma}

\begin{proof}
Use braid generators satisfying
\[
 (g_i-q)(g_i+1)=0,\qquad
 g_i^2=(q-1)g_i+q.
\]
For \(w\in S_4\), let \(g_w\) be a reduced product and
\(\ell(w)\) its length. The two idempotents are
\[
 e_+
 =\frac{\sum_{w\in S_4}g_w}
        {\sum_{w\in S_4}q^{\ell(w)}},
\qquad
 e_-
 =\frac{\sum_{w\in S_4}(-q^{-1})^{\ell(w)}g_w}
        {\sum_{w\in S_4}q^{-\ell(w)}}.
\]
Their denominators are nonzero at the sixth root used here. Put
\[
 z=\mu_\eta(g_i)=q-(1+q)\eta.
\]
Reducing the \(24\) terms by the quadratic relation, cyclicity, and the
Markov rule gives
\[
 \mu_\eta(e_+)
 =
 \frac{(z+1)((q+1)z+1)((q^2+q+1)z+1)}
 {(q+1)^2(q^2+1)(q^2+q+1)},
\]
\[
 \mu_\eta(e_-)
 =
 -\frac{(z-q)((q+1)z-q^2)
 ((q^2+q+1)z-q^3)}
 {(q+1)^2(q^2+1)(q^2+q+1)}.
\]
Substitute \(z=q-(1+q)\eta\) and use \(q^2-q+1=0\). This gives
\eqref{eq:four-strand-traces}. The common-zero assertion is immediate
from the displayed factors. The finite reduction is also replayed exactly
by the accompanying verifier.
\end{proof}

\begin{theorem}[Square-restriction inheritance]
\label{thm:square-inheritance}
If \(0\ne W\subseteq V\) is square-invariant, then \(\dim W\) is even and
the restricted operator satisfies
\[
  R|_{W\otimes W}\in
  \left[e^{i\pi/3},\frac12,\dim W\right].
\]
In particular, the restriction is non-scalar.
\end{theorem}

\begin{proof}
Square invariance makes \(W^{\otimes n}\) invariant under every local
generator, at every \(n\). Proposition~\ref{prop:trace-quotient} and
Lemma~\ref{lem:four-strand-traces} show that the ambient representation
kills both \(e_+\) and \(e_-\), because their images are orthogonal
projections of trace zero. Their restrictions to \(W^{\otimes4}\) also
vanish.

Put \(r=\dim W\) and let \(\eta_W\) be the normalized negative spectral
rank of \(R|_{W\otimes W}\), allowing temporarily the scalar values zero
and one. The restricted unitary again has no opposite eigenvalue pair, so
its normalized tensor trace is Markov. Applying
\eqref{eq:four-strand-traces} to the two zero restricted idempotents makes
\(\eta_W\) a common zero of the displayed polynomials. Hence
\(\eta_W=1/2\). Its negative spectral rank is \(r^2/2\), so \(r\) is even.
\end{proof}

\begin{corollary}[Restrictable descent]\label{cor:restrictable}
If a solution in dimension \(d\equiv2\pmod4\) is restrictable, its two
diagonal restrictions are balanced exceptional solutions of smaller even
dimensions, exactly one of which is again congruent to \(2\pmod4\).
Consequently a least-dimensional solution in the unresolved congruence
class is nonrestrictable, and no dimension-six solution is restrictable.
\end{corollary}

\begin{proof}
Apply Theorem~\ref{thm:square-inheritance} to \(W\) and \(W^\perp\).
Their even dimensions sum to \(2\bmod4\), so exactly one is
\(2\bmod4\). For \(d=6\), a proper even split is \(2+4\), but the
dimension-two exceptional class is empty~\cite{Lechner2026}.
\end{proof}

\begin{remark}[The one-sided gap]\label{rem:one-sided-gap}
Theorem~\ref{thm:square-inheritance} rules out a two-dimensional
square-invariant subspace in \(d=6\), since it would inherit the empty
dimension-two class. It does not rule out a four-dimensional
square-invariant \(W\). Corollary~\ref{cor:restrictable} says only that,
if such \(W\) exists, \((W^\perp)^{\otimes2}\) is not invariant.
\end{remark}

We now identify the precise common-leg algebra in dimension six.

\begin{lemma}[No common rank-two projection]\label{lem:no-common-rank2}
If \(d=6\), there is no rank-two projection in
\(\cC_L(P)\cap\cC_R(P)\).
\end{lemma}

\begin{proof}
Suppose \(z\) is such a projection and put \(W=zV\), so \(\dim W=2\).
Both tensor legs reduce \(H\), and
\[
  K=H|_{W\otimes W}
\]
is a Hermitian involution satisfying the exceptional cubic on
\(W^{\otimes3}\). Let \(Q=(\Id-K)/2\) and \(r=\rank Q\).

Ranks \(r=1,3\) are impossible by a short determinant gap. In the
rank-one case write
\[
 Q=|\operatorname{vec}S\rangle
   \langle\operatorname{vec}S|,
 \qquad \Tr(S^*S)=1.
\]
On \(\operatorname{ran}Q_{12}\), the two-projection relation makes the
compression \(Q_{12}Q_{23}Q_{12}\) have eigenvalues in
\(\{1,1/3\}\). Its determinant is therefore at least \(1/9\).
Vectorization computes the same determinant as \(|\det S|^4\), while
the arithmetic--geometric mean inequality gives
\[
 |\det S|^4\le\frac1{16}<\frac19.
\]
Complementation excludes rank three. Rank two would be a member of the
empty balanced dimension-two class. Hence \(r=0\) or \(4\), and
\[
  K=\varepsilon\Id_{W\otimes W},
  \qquad\varepsilon\in\{\pm1\}.
\]

Now restrict the cubic to \(W\otimes W\otimes V\). This space is invariant
because \(z\) lies in both leg commutants. The first adjacent copy is
\(X=\varepsilon\Id\); call the second \(Y\). Since \(X\) is scalar and
\(Y^2=\Id\),
\[
 XYX-YXY=Y-X.
\]
The exceptional cubic instead makes this \((X-Y)/3\), so \(Y=X\).
Therefore
\[
  H|_{W\otimes V}=\varepsilon\Id_{W\otimes V}.
\]
Taking the second partial trace contradicts
Corollary~\ref{cor:zero-marginals}:
\[
  0=z(\Tr_2H)z=6\varepsilon\Id_W.
\]
\end{proof}

\begin{theorem}[Scalar intersection of the two leg commutants]
\label{thm:common-leg}
Every dimension-six exceptional projection satisfies
\[
  \boxed{\cC_L(P)\cap\cC_R(P)=\bbC\Id_6.}
\]
\end{theorem}

\begin{proof}
A non-scalar finite-dimensional \(C^*\)-subalgebra of \(M_6\) contains a
nontrivial projection. Theorem~\ref{thm:leg-projection} makes its rank
even, hence two or four. In the latter case the complement has rank two.
Both alternatives contradict Lemma~\ref{lem:no-common-rank2}.
\end{proof}

\begin{remark}
Theorem~\ref{thm:common-leg} concerns the intersection after identifying
the two local copies of \(M_6\). It does not prove
\(\cC_L(P)=\bbC\Id_6\) or \(\cC_R(P)=\bbC\Id_6\) separately.
Exact relative-position models show that two non-scalar even-atom
subalgebras of \(M_6\) may have scalar intersection, so this distinction
is genuine.
\end{remark}

\section{The unrestricted Schmidt-rank-four reduction}
\label{sec:osr4}

Let \(H\) be an exceptional reflection with \(\OSR(H)=4\). Its intrinsic
left and right Schmidt supports are the slice spaces
\[
\begin{split}
 \cA&=\{(\operatorname{id}\otimes\omega)(H):
                  \omega\in M_d^*\},\\
 \cB&=\{(\omega\otimes\operatorname{id})(H):
                  \omega\in M_d^*\}.
\end{split}
\]
They are four-dimensional and independent of a chosen Schmidt
decomposition. Put
\[
  \cO_{\cA}=\bbC\Id+\cA,\qquad
  \cO_{\cB}=\bbC\Id+\cB.
\]
For \(T\in M_d\), write \([T]_{\cA}\) and \([T]_{\cB}\) for its
cosets in \(M_d/\cO_{\cA}\) and \(M_d/\cO_{\cB}\), respectively.

For subspaces \(\cX,\cY\subseteq M_d\), define the intrinsic
\emph{joint-sandwich map}
\begin{equation}
 \cS_{\cY\mid\cX}:
 \cY\otimes\cY\longrightarrow\operatorname{Hom}_{\bbC}(\cX,M_d),
 \qquad
 \sum_\ell y_\ell\otimes z_\ell
 \longmapsto
 \left[x\mapsto\sum_\ell y_\ell xz_\ell\right].
 \label{eq:sandwich-map}
\end{equation}

Because \(H\) is Hermitian, a real singular-value decomposition in the
real Hilbert spaces of Hermitian matrices gives
\begin{equation}
 H=\sum_{i=1}^4\sigma_iA_i\otimes B_i,
 \qquad \sigma_i>0,
 \label{eq:hermitian-schmidt}
\end{equation}
where \(A_i,B_i\) are Hermitian and Hilbert--Schmidt orthonormal.
Corollary~\ref{cor:zero-marginals} makes every \(A_i,B_i\) traceless.
Thus \(\cO_{\cA},\cO_{\cB}\) are five-dimensional operator systems.

\begin{lemma}[Full-cubic quotient identities]\label{lem:quotient-identities}
For every \(x\in\cO_{\cA}\) and \(y\in\cO_{\cB}\),
\begin{equation}
 \sum_{i,k}\sigma_i\sigma_k
 [A_iA_k]_{\cA}\otimes B_i xB_k=0
 \quad\text{in }
 (M_d/\cO_{\cA})\otimes M_d,
 \label{eq:quotient-A}
\end{equation}
\begin{equation}
 \sum_{i,k}\sigma_i\sigma_k
 A_i yA_k\otimes[B_iB_k]_{\cB}=0
 \quad\text{in }
 M_d\otimes(M_d/\cO_{\cB}).
 \label{eq:quotient-B}
\end{equation}
\end{lemma}

\begin{proof}
It is clearest to begin with an arbitrary minimal factorization
\[
  H=\sum_{i=1}^4a_i\otimes b_i.
\]
The two cubic words expand as
\[
\begin{split}
 H_{12}H_{23}H_{12}
 &=\sum_{i,j,k}a_ia_k\otimes b_ia_jb_k\otimes b_j,\\
 H_{23}H_{12}H_{23}
 &=\sum_{i,j,k}a_j\otimes a_ib_ja_k\otimes b_ib_k.
\end{split}
\]
Project the first tensor leg of~\eqref{eq:cubic-H} to
\(M_d/(\bbC\Id+\cA)\). The second cubic word and both linear terms
vanish there. Applying an arbitrary functional to the third leg, and
using that the \(b_j\) and \(a_j\) are bases of the two Schmidt supports,
gives
\[
  \sum_{i,k}[a_ia_k]_{\cA}\otimes b_i x b_k=0
  \qquad(x\in\cA).
\]
Projecting the first leg of \(H^2=\Id\) supplies the same identity at
\(x=\Id\). This proves the all-\(x\) statement on
\(\cO_{\cA}\). Tensor reversal gives the symmetric identity.
Substituting~\eqref{eq:hermitian-schmidt} yields
\eqref{eq:quotient-A}--\eqref{eq:quotient-B}.
\end{proof}

\begin{theorem}[Joint-sandwich divisibility]\label{thm:sandwich}
If either
\[
  \cS_{\cB\mid\cO_{\cA}}
  \quad\text{or}\quad
  \cS_{\cA\mid\cO_{\cB}}
\]
is injective, then \(4\mid d\).
\end{theorem}

\begin{proof}
Suppose \(\cS_{\cB\mid\cO_{\cA}}\) is injective. Apply a linear
functional \(\varphi\) on \(M_d/\cO_{\cA}\) to the first factor of
\eqref{eq:quotient-A}. The coefficient matrix
\[
  c_{ik}=\sigma_i\sigma_k
  \varphi([A_iA_k]_{\cA})
\]
satisfies
\[
  \sum_{i,k}c_{ik}B_i xB_k=0
  \qquad(x\in\cO_{\cA}).
\]
Injectivity makes every \(c_{ik}\) zero. Since \(\varphi\) was arbitrary,
\[
  A_iA_k\in\bbC\Id+\cA
  \quad\text{for all }i,k.
\]
Therefore the five-dimensional operator system
\(\bbC\Id+\cA\) is closed under multiplication. It is a finite-dimensional
\(C^*\)-subalgebra of \(M_d\).

The only complex \(C^*\)-algebras of dimension five are
\[
  \bbC^5,\qquad M_2(\bbC)\oplus\bbC,
\]
because their dimensions are sums of squares. Every faithful unital
concrete representation of either algebra has a rank-one projection in
its commutant: take a rank-one subprojection in a nonzero multiplicity
space, using the nonzero scalar summand in the second case.

Independence of the \(B_i\) gives
\[
 [X\otimes\Id,H]=0
 \quad\Longleftrightarrow\quad
 [X,A_i]=0\ \text{for every }i.
\]
Thus this rank-one projection lies in \(\cC_L(P)\).
Theorem~\ref{thm:leg-projection} forces \(8\mid d^2\), hence \(4\mid d\).
The other sandwich map is handled with the legs exchanged.
\end{proof}

\begin{corollary}[Simultaneous rank-four degeneracy]
\label{cor:sandwich-degeneracy}
Suppose \(d\equiv2\pmod4\) and \(\OSR(H)=4\). Then both maps in
Theorem~\ref{thm:sandwich} are noninjective. More precisely, in the
Hermitian Schmidt gauge~\eqref{eq:hermitian-schmidt}, each kernel contains
a nonzero tensor whose \(4\times4\) coefficient matrix is Hermitian and
traceless.
\end{corollary}

\begin{proof}
Neither \(\cO_{\cA}\) nor \(\cO_{\cB}\) can be an algebra, since the
five-dimensional \(C^*\)-algebra argument in the preceding proof would
again give a rank-one leg projection and contradict
\(d\equiv2\pmod4\).

The quotient-product span generated by
\([A_iA_k]_{\cA}\) is therefore nonzero and is closed under adjoint.
Choose a nonzero functional \(\varphi\) on this span satisfying
\(\varphi(z^*)=\overline{\varphi(z)}\). Then
\[
  C_{ik}
  =\sigma_i\sigma_k\varphi([A_iA_k]_{\cA})
\]
is a nonzero Hermitian coefficient matrix, and
\eqref{eq:quotient-A} says
\[
  \sum_{i,k}C_{ik}B_i xB_k=0
  \quad(x\in\cO_{\cA}).
\]
At \(x=\Id\), take ordinary trace. Hilbert--Schmidt orthonormality of the
\(B_i\) gives
\[
  0=\sum_{i,k}C_{ik}\Tr(B_iB_k)=\Tr C.
\]
The symmetric argument gives the annihilator on the other leg.
\end{proof}

\begin{corollary}[A low-Kraus-rank boundary map]\label{cor:cp-boundary}
Under the hypotheses of Corollary~\ref{cor:sandwich-degeneracy}, there is
on each leg a unital completely positive map of Kraus rank at most three
which acts as \(-1/3\) on the opposite four-dimensional Schmidt support.
\end{corollary}

\begin{proof}
Let \(D=\operatorname{diag}(\sigma_1^2,\ldots,\sigma_4^2)>0\), and let
\(C=C^*\ne0\) be a traceless annihilator from the preceding proof.
Move from \(D\) along the line \(D+tC\) to a nonzero boundary point
\[
  G=D+t_*C\ge0,\qquad 1\le\rank G\le3.
\]
The endpoint cannot be zero, since that would make \(C\) proportional to
the positive matrix \(D\), contradicting \(\Tr C=0\). Define
\[
  \Psi_B(x)
  =\frac1d\sum_{i,k}G_{ik}B_i xB_k.
\]
Positivity of \(G\) makes this completely positive with Kraus rank at most
three. The contraction of \(H^2=\Id\) gives
\[
  \frac1d\sum_i\sigma_i^2B_i^2=\Id,
\]
while a partial trace of the cubic gives
\[
  \frac1d\sum_i\sigma_i^2B_iA_jB_i=-\frac13A_j.
\]
The \(C\)-direction vanishes on \(\bbC\Id+\cA\), so \(\Psi_B\) has the
same two actions. The other leg is identical.
\end{proof}

\begin{remark}
The map in Corollary~\ref{cor:cp-boundary} is not automatically trace
preserving. That would require the transpose coefficient direction to
annihilate the identity as well. No such transpose invariance is assumed
or proved.
\end{remark}

\section{Leg-commutant arithmetic and low Schmidt rank}
\label{sec:low-osr}

We first isolate the arithmetic consequence of a genuine one-leg
projection. This is the bridge from tensor-local structure to
four-divisibility.

\begin{lemma}[Two exceptional projections]\label{lem:two-projections}
Let \(a,b\) be projections on a \(D\)-dimensional Hilbert space such that
\[
  aba-bab=\frac13(a-b),
\]
\[
  \rank a=\rank b=\frac D2,\qquad
  \Tr(ab)=\frac D4.
\]
Then
\[
 \dim(\operatorname{ran}a\cap\operatorname{ran}b)
 =
 \dim(\ker a\cap\ker b)
 =\frac D8.
\]
\end{lemma}

\begin{proof}
Put
\[
  z=aba-\frac13a=bab-\frac13b.
\]
Projection identities give \(az=za=z\) and \(bz=zb=z\). Thus the range
of \(z\) lies in the common range of \(a,b\), where \(z=2\Id/3\).
Self-adjointness therefore gives \(z=(2/3)e\), with \(e\) the
common-range projection. Taking traces,
\[
  \Tr e
  =\frac{\Tr(ab)-\frac13\Tr a}{2/3}
  =\frac D8.
\]
Equal half ranks make the common-range and common-kernel dimensions equal;
this also follows by applying the same argument to \(\Id-a,\Id-b\).
\end{proof}

\begin{theorem}[Invariant leg projection]\label{thm:leg-projection}
If a rank-\(r\) projection \(z\) belongs to either
\(\cC_L(P)\) or \(\cC_R(P)\), then
\[
  \boxed{8\mid rd^2.}
  \label{eq:leg-divisibility}
\]
\end{theorem}

\begin{proof}
Suppose \(z\in\cC_L(P)\), put \(W=zV\), and let \(Q\) be the
restriction of \(P\) to \(W\otimes V\). Theorem~\ref{thm:standardness}
gives
\[
  \rank Q
  =\Tr((z\otimes\Id)P)
  =\Tr\!\left(z\,\Tr_2P\right)
  =\frac{rd}{2}.
\]
Restrict \(P_1,P_2\) to \(W\otimes V\otimes V\). The resulting
projections \(a=Q\otimes\Id\), \(b=\Id_W\otimes P\) act on a space of
dimension \(D=rd^2\), obey the exceptional relation, and both have rank
\(D/2\). Their overlap trace is
\[
\begin{aligned}
 \Tr(ab)
 &=\Tr_{W\otimes V}\!
   \left[Q(\Id_W\otimes\Tr_2P)\right]\\
 &=\frac d2\Tr Q=\frac{rd^2}{4}.
\end{aligned}
\]
Lemma~\ref{lem:two-projections} makes their common-range dimension
\(rd^2/8\), an integer. The right-leg case follows after tensor flip.
\end{proof}

\begin{corollary}\label{cor:odd-leg}
If \(d\equiv2\pmod4\), every projection in either leg commutant has even
rank. In particular, either leg containing a rank-one projection forces
\(4\mid d\).
\end{corollary}

We next show that low operator-Schmidt rank supplies such a projection.
The subtle point is that the known rank-two and rank-three unitary
structure theorems use four independent local pre- and post-unitaries.
That equivalence does not preserve the Yang--Baxter equation.

\begin{theorem}[Twisted control]\label{thm:twisted-control}
Let \(H\in\End(V\otimes V)\) be a Hermitian unitary which is locally
equivalent, through independent product unitaries before and after, to a
controlled bipartite unitary. Suppose
\[
  H_1H_2H_1-H_2H_1H_2=c(H_1-H_2)
  \label{eq:general-cubic}
\]
with \(|c|\ne1\). Then one true one-leg commutant of \(H\) contains a
rank-one projection.
\end{theorem}

\begin{proof}
Assume first that the supplied control is on the first factor. There are
unitaries \(Q,R,S,T\), rank-one coordinate projections \(E_i\), and
unitaries \(U_i\) such that
\[
  H=(Q\otimes S)
    \left(\sum_iE_i\otimes U_i\right)
    (R\otimes T).
  \label{eq:four-local}
\]
We do not insert the controlled operator in parentheses into
\eqref{eq:general-cubic}. Instead, apply the valid same-site conjugacy
\[
  \widetilde H=(Q^*\otimes Q^*)H(Q\otimes Q).
\]
It preserves~\eqref{eq:general-cubic} and gives the twisted controlled
form
\begin{equation}
  \widetilde H=\sum_iE_iK\otimes V_i,
  \qquad K=RQ,\quad V_i=Q^*SU_iTQ.
  \label{eq:twisted-form}
\end{equation}

The \(ij\) block of \(\widetilde H\) is \(K_{ij}V_i\). Hermiticity is
therefore
\begin{equation}
  K_{ij}V_i=\overline{K_{ji}}V_j^*.
  \label{eq:hermitian-block}
\end{equation}
In particular, the support graph of \(K\) is undirected, and along each
edge the projective classes of the target unitaries alternate between
\([W]\) and \([W^*]\).

If a connected component of this graph is nonbipartite, an odd cycle
identifies these two classes. On that component,
\[
  \widetilde H_C=A_C\otimes W
\]
for a unitary \(A_C\). A rank-one spectral projection of the normal
matrix \(A_C\), extended by zero across the other graph components,
commutes with the full \(\widetilde H\) on its first leg.

Suppose instead that every component is bipartite. Unitarity of \(K\)
shows that the two parts in each component have equal dimension.
Absorbing phases and choosing matched orthonormal bases puts
\(\widetilde H\) in the fixed-point-free form
\begin{equation}
  \widetilde H
  =\sum_x|\bar x\rangle\langle x|\otimes U_x,
  \qquad U_{\bar x}=U_x^*,
  \label{eq:fixed-point-free}
\end{equation}
where \(x\mapsto\bar x\) is an involution without fixed points.

Apply the first-leg functional
\(\varphi_x(A)=\langle\bar x|A|x\rangle\) to
\eqref{eq:general-cubic}. Products of two first-leg coefficients in
\eqref{eq:fixed-point-free} are diagonal whenever nonzero. Hence the
\(\widetilde H_1\widetilde H_2\widetilde H_1\) coefficient vanishes,
whereas the other cubic word contributes
\[
  \widetilde H(U_x\otimes\Id)\widetilde H.
\]
The two linear terms contribute \(U_x\otimes\Id\) and zero. We obtain
\[
  -\widetilde H(U_x\otimes\Id)\widetilde H
  =c(U_x\otimes\Id),
\]
which equates a unitary with an operator of norm \(|c|\ne1\), a
contradiction. Thus a nonbipartite component exists and yields the claimed
rank-one projection. Conjugating back preserves its rank.

If the original controlled side is the second factor, tensor flip
interchanges the two legs and preserves~\eqref{eq:general-cubic}; the same
argument applies.
\end{proof}

Cohen--Yu prove that every bipartite unitary of operator-Schmidt rank two
is locally equivalent to a controlled unitary~\cite[Theorem 6]{CohenYu}.
Chen--Yu prove the analogous statement in rank three
\cite[Theorem 11]{ChenYu}. Their hypotheses allow arbitrary finite
bipartite dimensions and the four-local equivalence used in
\eqref{eq:four-local}.

\begin{theorem}[Complete low-Schmidt spectrum]\label{thm:low-osr}
There exists an exceptional reflection with \(\OSR(H)\le3\) in local
dimension \(d\) if and only if
\[
  \boxed{4\mid d.}
\]
\end{theorem}

\begin{proof}
For necessity, apply Theorem~\ref{thm:twisted-control} with \(c=1/3\)
and then Corollary~\ref{cor:odd-leg}. Rank one is immediate, and ranks
two and three are covered by the cited controlled-unitary theorems.

For sufficiency, the published \(d=4\) reflection has
operator-Schmidt rank three~\cite{Kriebel2026d4}. If \(d=4m\), identity
stabilization replaces each Schmidt factor \(A_j,B_j\) by
\(A_j\otimes\Id_m,B_j\otimes\Id_m\), preserving both the exceptional
relations and Schmidt rank three.
\end{proof}

\begin{remark}
Theorem~\ref{thm:low-osr} does not claim that four-local equivalence
preserves Yang--Baxter locality. The only such transformation used in the
proof is the same-site conjugacy leading to~\eqref{eq:twisted-form}.
Hermiticity and the cubic relation then provide the new argument.
\end{remark}

\section{The surviving dimension-six branch}
\label{sec:surviving-branch}

Combining the preceding results gives the promised exact frontier.

\begin{corollary}[Exact dimension-six frontier]
\label{cor:d6-frontier}
Let \(P\) be an exceptional projection on
\(\bbC^6\otimes\bbC^6\), let \(H=\Id-2P\), and let \(R\) be given
by~\eqref{eq:H-R}. Then:
\begin{enumerate}[label=\textup{(\roman*)}]
\item
\[
  \OSR(H)\ge4,\qquad \OSR(R)\ge5.
\]
If \(\OSR(H)=4\), both intrinsic joint-sandwich maps in
Theorem~\ref{thm:sandwich} are singular and have nonzero Hermitian
traceless coefficient-matrix annihilators.
\item \(R\) is nonrestrictable. It has no two-dimensional
square-invariant local subspace.
\item
\[
  \cC_L(P)\cap\cC_R(P)=\bbC\Id_6.
\]
\item \(H\) is not diagonal in a primitive generalized Bell basis and is
not a four-product Clifford-frame reflection of
Theorem~\ref{thm:clifford}.
\end{enumerate}
\end{corollary}

\begin{proof}
Part (i) is Theorem~\ref{thm:low-osr},
\eqref{eq:osr-shift}, and
Corollary~\ref{cor:sandwich-degeneracy}. Part (ii) is
Theorem~\ref{thm:square-inheritance} and
Corollary~\ref{cor:restrictable}. Part (iii) is
Theorem~\ref{thm:common-leg}. Part (iv) follows from
Theorems~\ref{thm:bell} and~\ref{thm:clifford}.
\end{proof}

\begin{remark}[What ``primitive'' means here]
\label{rem:primitive-word}
Corollary~\ref{cor:d6-frontier} does not assert irreducibility in every
tensor-local sense. It leaves open a four-dimensional square-invariant
subspace \(W\), provided that the complementary square
\((W^\perp)^{\otimes2}\) is not invariant. It also leaves open non-scalar
individual algebras \(\cC_L(P)\) and \(\cC_R(P)\) with scalar
intersection. Thus the safe summary is that a dimension-six solution
would be \emph{nonrestrictable} and would have no common one-leg
symmetry, in the precise senses of
Definition~\ref{def:restrictable} and~\eqref{eq:leg-commutants}.
\end{remark}

\subsection{The bounded one-sided
\texorpdfstring{\(4+2\)}{4+2} reduction}

The remaining square-restriction possibility admits a finite exact block
formulation. Write \(V=W\oplus U\), with
\(\dim W=4\), \(\dim U=2\), and assume only that \(W^{\otimes2}\)
reduces \(H\). By Theorem~\ref{thm:square-inheritance}, the diagonal
\(W^{\otimes2}\) block is a balanced \(d=4\) exceptional reflection.
Coloring the three tensor positions by \(W,U\) decomposes the cubic into
all \(2^6=64\) input-output color blocks. Unitarity and the zero-marginal
identities add exact quadratic and linear equations for the remaining
blocks.

This reduction has been written out and replayed symbolically in the
supplement. It does \emph{not} force \(U^{\otimes2}\) to be invariant:
the coefficient system still contains genuine leakage blocks. Frozen-core
numerical searches found no solution, but those residuals carry no
nonexistence content and are not used in any theorem. This is the precise
point at which the bounded one-sided attack stopped.

\subsection{A secondary four-dimensional sitewise orbit}

The exact computations also produce a circle of \(d=4\) reflections in a
three-term color/face form. Exact conjugation shows that this entire circle
is one orbit under sitewise conjugacy \(H\mapsto(U\otimes U)H
(U\otimes U)^*\). The invariant
\[
  \Tr((HF)^4),
\]
where \(F\) is tensor flip, equals \(-16/3\) on this orbit and \(16\) on
the published five-Pauli witness. The two exhibited matrices are therefore
not sitewise-unitarily conjugate. They are nevertheless equivalent under
Lechner's broader braid-character relation, since they have the same
triple \([q,\eta,d]\). This secondary result is documented in the
supplement; it is not a classification of all \(d=4\) sitewise orbits and
no separate novelty claim is made for it.

\section{Dimension table and open problem}
\label{sec:open-problem}

The exact status after this work is:
\[
\begin{array}{c|c|l}
\text{local dimension}&\text{status}&\text{reason}\\ \hline
d\text{ odd}&\text{empty}&d^2/2\notin\mathbb Z\\
2&\text{empty}&\text{Lechner's dimension-two exclusion}\\
4m,\ m\ge1&\text{nonempty}&d=4\text{ witness and identity stabilization}\\
4m+2,\ m\ge1&\text{open}&\text{no universal obstruction or exact witness}
\end{array}
\]

\begin{problem}[Exceptional dimension spectrum]
\label{prob:spectrum}
Determine whether
\[
 \left[e^{i\pi/3},\frac12,d\right]\ne\varnothing
 \quad\Longleftrightarrow\quad 4\mid d,
\]
or construct an exact solution in some dimension
\(d\equiv2\pmod4\), beginning with \(d=6\).
\end{problem}

Theorem~\ref{thm:tower-summary} says that the abstract Jones--Wenzl tower
cannot decide Problem~\ref{prob:spectrum} through simple multiplicities,
central ranks, Markov weights, or their represented branching recurrences.
Corollary~\ref{cor:d6-frontier} identifies what a first counterexample must
avoid. The rank-four sandwich identities are the only result here which
applies to unrestricted Schmidt supports at the first surviving rank; they
leave open the simultaneous-degeneracy branch.

\section{Exact verification and provenance}
\label{sec:verification}

All theorem statements have human-readable proofs above. Computation is
used only to replay finite algebraic reductions, calibrate formulas, and
document limitation models. The public repository contains:
\begin{enumerate}[label=\textup{(\roman*)}]
\item a theorem-dependency graph separating imported results, proved
lemmas, and exact computational certificates;
\item a verifier manifest specifying the claim, arithmetic domain, and
scope of each program;
\item a one-command exact verifier suite for the central results;
\item a supplementary catalogue separating exact ansatz theorems from
numerical evidence and failed construction mechanisms;
\item seeds, commands, dependency information, raw outputs, and claim
statuses for every new numerical experiment.
\end{enumerate}

The central exact programs independently replay the Hecke multiplicities,
two-projection block arithmetic, invariant-leg divisibility, the
low-Schmidt contraction identities, the unrestricted rank-four quotient
identities, the four-strand restriction polynomials, the dimension-six
commutant block analysis, primitive-Weyl coefficient formulas, and the
four-product Clifford parity calculation. They are checks of finite
identities, not substitutes for the logical implications proved here.

The original numerical search which discovered the published \(d=4\)
five-Pauli support was not preserved and is not described as reproducible.
Every computation performed for the present paper has a retained source,
command, seed when applicable, raw result, and commit history.

\section*{Priority and limitations}

Targeted searches covered unitary Hecke tensor-space representations,
controlled bipartite unitaries of Schmidt ranks two and three, generalized
Yang--Baxter operators, quaternionic and Clifford realizations, Ocneanu
cell systems, and the operator-algebraic theory of Yang--Baxter
endomorphisms. We found no prior statement of the automatic-standardness
consequence in this exceptional matrix class, the exact all-level
multiplicity limitation, the twisted-control four-divisibility theorem, or
the joint-sandwich rank-four obstruction. These results should therefore
be described as \emph{apparently new} pending specialist review.

The paper does not settle Problem~\ref{prob:spectrum}, exclude the
simultaneously singular rank-four branch, exclude a one-sided
four-dimensional square restriction in \(d=6\), classify all
operator-Schmidt-rank-four unitaries, or classify all \(d=4\) solutions.
Negative numerical searches and the supplementary finite ansatz
exclusions are not evidence of global nonexistence.

\section*{AI-assistance disclosure}

OpenAI language models were used extensively to propose proof routes,
derive and check symbolic identities, write and review code, search the
literature, draft prose, and conduct adversarial audits. The author directed
the research program and accepts responsibility for every claim. Central
results were rederived in human-readable form and, where finite algebra was
involved, replayed by exact programs with retained outputs. No model
residual, confidence score, or unsupported model assertion is treated as
mathematical evidence. The complete work trace, including failed
approaches and proof repairs, is published with the paper.

\begin{thebibliography}{99}

\bibitem{ChenYu}
L.~Chen and L.~Yu,
\newblock On the Schmidt-rank-three bipartite and multipartite unitary
operator,
\newblock \emph{Ann. Physics} \textbf{351} (2014), 682--703.
\newblock \url{https://doi.org/10.1016/j.aop.2014.09.026}

\bibitem{CohenYu}
S.~M. Cohen and L.~Yu,
\newblock All unitaries having operator Schmidt rank \(2\) are controlled
unitaries,
\newblock \emph{Phys. Rev. A} \textbf{87} (2013), 022329.
\newblock \url{https://doi.org/10.1103/PhysRevA.87.022329}

\bibitem{ContiLechner}
R.~Conti and G.~Lechner,
\newblock Yang--Baxter endomorphisms,
\newblock \emph{J. London Math. Soc.} \textbf{103} (2021), 633--671.
\newblock \url{https://doi.org/10.1112/jlms.12387}

\bibitem{GHR2013}
C.~Galindo, S.-M.~Hong, and E.~C.~Rowell,
\newblock Generalized and quasi-localizations of braid group
representations,
\newblock \emph{Int. Math. Res. Not. IMRN} \textbf{2013}, no.~3,
693--731.
\newblock \url{https://doi.org/10.1093/imrn/rnr269}

\bibitem{Kriebel2026d4}
A.~Kriebel,
\newblock A minimal unitary localization of the
\(\cC(\mathfrak{sl}_3,6)\) Jones--Wenzl representations,
\newblock unrefereed research note, 2026.
\newblock
\url{https://aleckriebel.github.io/Math/papers/exceptional-ybe-d4/}

\bibitem{Lechner2026}
G.~Lechner,
\newblock The classification problem for unitary \(R\)-matrices with two
eigenvalues,
\newblock arXiv:2603.20158v1, 2026.
\newblock \url{https://arxiv.org/abs/2603.20158}

\bibitem{Wenzl1988}
H.~Wenzl,
\newblock Hecke algebras of type \(A_n\) and subfactors,
\newblock \emph{Invent. Math.} \textbf{92} (1988), 349--383.
\newblock \url{https://doi.org/10.1007/BF01404457}

\end{thebibliography}

\end{document}
